\begin{frame}{CTR extraction}
\small
\begin{columns}[T,totalwidth=\textwidth]
\begin{column}{0.58\textwidth}
\begin{block}{Direct histogram FWHM}
The coincidence-time residuals are histogrammed with a fixed physical bin width
\[
\Delta t_{\rm bin}=5~\mathrm{ps}\quad\text{(default, configurable).}
\]
Let $N_{\max}$ be the highest histogram count. The left and right crossings of $N_{\max}/2$ are obtained by linear interpolation between adjacent bin centres.
\[
\CTR=t_{\mathrm{right},1/2}-t_{\mathrm{left},1/2}.
\]
\end{block}

\begin{block}{Statistical uncertainty}
The selected events are resampled with replacement and the same fixed-bin FWHM estimator is repeated. The default uncertainty uses 100 bootstrap replicas:
\[
\sigma_{\CTR}=\operatorname{std}(\CTR_1,\ldots,\CTR_{100}).
\]
\end{block}
\end{column}

\begin{column}{0.40\textwidth}
\begin{block}{Timing difference}
For each selected coincidence event:
\[
\Delta t=t_1-t_2.
\]
The same estimator is applied to LED, CFD and ML-corrected residuals.
\end{block}

\begin{block}{Why direct FWHM?}
The metric measures the width of the central timing peak directly and does not require a parametric model for the complete residual distribution.
\end{block}
\end{column}
\end{columns}
\mainmsg{CTR is the direct FWHM of a fixed-bin timing histogram; reporting plots may use different display binning.}
\end{frame}
